% Balance vs. imbalance: an empirical refutation of the "avoid unbalanced
% splits" heuristic. Data generated 2026-06 from the catalog's best
% non-commutative R-ranks (snapshot date in the table captions); the driver
% lives in the search package (AllocationOptimizer + the omega/imbalance probe).

A recurring heuristic --- explicit in DIS09 and folklore well before --- is that
recombination should prefer \emph{balanced} block splits and avoid unbalanced
ones. We test this directly against the catalog and reach a sharper, two-sided
conclusion: \textbf{balance is better on average, but the records are
unbalanced.} A search that prunes unbalanced shapes is therefore average-case
sound yet \emph{record-blind} --- it discards exactly the shapes that set the
lowest exponents.

Throughout this section the per-shape exponent is
\(\omega(n,m,p) = 3\,\ln R(n,m,p) / \ln(nmp)\), where \(R\) is the best
catalogued non-commutative rank over \(\Reals\) (catalog snapshot
2026-06; the catalog holds a single \Reals-only scheme, so the same
figures hold over \(\Rationals\), the field all later sweeps default
to).

\subsection{On average, \texorpdfstring{\(\omega\)}{omega} rises with imbalance}
\label{ssec:omega-imbalance}

Table~\ref{tab:omega-imbalance} buckets every catalogued shape
\(\nmpshape{n}{m}{p}\) with \(2 \le n \le m \le p \le 16\) by its imbalance
\(p-n\), and reports the median \(\omega\) per bucket.

\begin{table}[h]
\centering
\caption{Median \(\omega\) over all catalogued shapes
\(2 \le n \le m \le p \le 16\), bucketed by imbalance \(p-n\). The median
climbs monotonically --- a \emph{typical} unbalanced shape is less efficient.}
\label{tab:omega-imbalance}
\begin{tabular}{rrr}
\toprule
imbalance \(p-n\) & \#shapes & median \(\omega\) \\
\midrule
0  & 15 & 2.8155 \\
1  & 28 & 2.8148 \\
2  & 39 & 2.8177 \\
3  & 48 & 2.8180 \\
4  & 55 & 2.8169 \\
5  & 60 & 2.8193 \\
6  & 63 & 2.8214 \\
7  & 64 & 2.8229 \\
8  & 63 & 2.8274 \\
9  & 60 & 2.8299 \\
10 & 55 & 2.8332 \\
11 & 48 & 2.8402 \\
12 & 39 & 2.8536 \\
13 & 28 & 2.8616 \\
14 & 15 & 2.8695 \\
\bottomrule
\end{tabular}
\end{table}

The median rises from \(2.8155\) (balanced) to \(2.8695\) at imbalance \(14\),
staying flat (\(\approx 2.817\)) out to imbalance \(\sim 7\) and then climbing
steeply --- the penalty for imbalance is convex, not linear. This is the
average-case fact the avoid-unbalanced heuristic rests on.

\subsection{\dots\ but the lowest exponents are unbalanced}
\label{ssec:omega-records}

The aggregate hides the tail. Among the same shapes, the fifteen smallest
\(\omega\) are almost all unbalanced, and the global minimiser in this range is
\(\nmpshape{3}{3}{6}\) at \(\omega = 2.7743\) --- below \emph{every} cubic:
\(\nmpshape{4}{4}{4}\) (\(2.7925\)), \(\nmpshape{8}{8}{8}\) (\(2.7974\)),
\(\nmpshape{12}{12}{12}\) (\(2.7957\)). The low-\(\omega\) frontier is dominated
by the \(2{:}1\)-ratio family \(\nmpshape{3}{3}{6}\), \(\nmpshape{6}{6}{12}\),
\(\nmpshape{9}{9}{12}\), \(\nmpshape{5}{5}{12}\), \(\nmpshape{3}{4}{6}\).
Pruning unbalanced splits removes precisely these record-holders.

\subsection{Why the advantage is not extractable by outer splitting}
\label{ssec:omega-not-extractable}

It is tempting to conclude that a recombination should \emph{seek} unbalanced
splits so its sub-blocks land on these efficient shapes. The catalog says
otherwise, and the mechanism is instructive. Consider \(\nmpshape{9}{9}{9}\)
recombined by Strassen \(\nmpshape{2}{2}{2}{=}7\): since \(9 = 5+4 = 6+3\), the
balanced split \((5,4)^3\) and the unbalanced \((6,3)^3\) both tile exactly (no
padding). Table~\ref{tab:9-faceoff} puts their seven-product multisets
side by side.

\begin{table}[h]
\centering
\caption{Strassen \(\nmpshape{2}{2}{2}\) recombination of
\(\nmpshape{9}{9}{9}\): the seven sub-products grouped by corner tier. The
unbalanced \((6,3)^3\) split saves \(63\) on the three one-max corners
(\(\nmpshape{3}{3}{6}{=}40\) vs \(61\)) but pays \(+60\) on the all-max corner
and \(+12\) on the middle tier --- net \(+9\), so balance wins here.}
\label{tab:9-faceoff}
\begin{tabular}{llrlrr}
\toprule
corner tier & \multicolumn{2}{c}{\((5,4)^3 \to 504\)} & \multicolumn{2}{c}{\((6,3)^3 \to 513\)} & \(\Delta\) \\
(multiplicity) & shape & subtotal & shape & subtotal & \\
\midrule
all-max \(\times 1\) & \(\nmpshape{5}{5}{5}{=}93\)  & 93  & \(\nmpshape{6}{6}{6}{=}153\) & 153 & \(+60\) \\
two-max \(\times 3\) & \(\nmpshape{5}{5}{4}{=}76\)  & 228 & \(\nmpshape{6}{6}{3}{=}80\)  & 240 & \(+12\) \\
one-max \(\times 3\) & \(\nmpshape{5}{4}{4}{=}61\)  & 183 & \(\nmpshape{3}{3}{6}{=}40\)  & 120 & \(-63\) \\
\midrule
total & & \textbf{504} & & \textbf{513} & \(+9\) \\
\bottomrule
\end{tabular}
\end{table}

A \(\nmpshape{2}{2}{2}\) split realises seven of the eight corners of
\(\{a_1,a_2\}\times\{b_1,b_2\}\times\{c_1,c_2\}\), and \emph{one of them is the
all-max corner} \((\max,\max,\max)\). Because \(R\) grows like
\(\text{size}^{\,\omega}\) with \(\omega \approx 2.8\), pushing the max block
from \(5\) to \(6\) costs more than the cheap corners recover. The efficiency
of \(\nmpshape{3}{3}{6}\) is real, but a single block split cannot isolate it
from its complementary all-max corner.

The corollary is constructive: an efficient unbalanced atom is exploited not by
\emph{splitting} a target down onto it, but by \emph{composing up} from it via
the Kronecker product, where the favourable ratio compounds instead of dragging
a complementary max-corner along (cf.~\ref{ssec:disjoint-sum}).

\subsection{Imbalance can win the split, too}
\label{ssec:imbalance-split-wins}

Section~\ref{ssec:omega-not-extractable} showed balance winning a single
example; it is not the rule. We swept every target
\(\nmpshape{n}{m}{p}\), \(3 \le n \le m \le p \le 16\), computing for a fixed
base the \emph{exact} rank-minimising allocation by the branch-and-bound of
Section~\ref{ssec:alloc-bnb} (balanced incumbent + admissible lower bound,
verified against exhaustive search),
and flagged every target whose optimal split is strictly cheaper than the
balanced split. The search is deliberately \emph{single-base} and
\emph{non-recursive}: sub-blocks are costed by the catalog's best rank as a
black box, so a ``win'' means the unbalanced split beats its own balanced
sibling for that base --- not that it beats the global record.

Unbalanced splits win for \(15\) targets under Strassen \(\nmpshape{2}{2}{2}\)
and \(120\) under Strassen--Winograd (whose asymmetric product supports reward
imbalance far more often). The effect is not confined to rectangular targets:
the cubic \(\nmpshape{13}{13}{13}\) is itself a win.
Table~\ref{tab:imbalance-wins} lists representatives.

\begin{table}[h]
\centering
\caption{Targets where the rank-optimal single-base split is \emph{unbalanced}
(strictly cheaper than the balanced split). Single base, no recursion;
sub-blocks costed by catalog SOTA. Savings are small but the existence is the
point: a hard ``always split balanced'' rule is provably suboptimal.}
\label{tab:imbalance-wins}
\begin{tabular}{llrrr}
\toprule
base & target & balanced & optimal split & best \\
\midrule
Strassen & \(\nmpshape{13}{13}{13}\) & 1434 & \((5,8)(6,7)(6,7)\) & 1432 \\
Strassen & \(\nmpshape{6}{6}{10}\)   & 252  & \((3,3)(3,3)(4,6)\) & 247 \\
Strassen & \(\nmpshape{12}{12}{14}\) & 1281 & \((6,6)(6,6)(6,8)\) & 1271 \\
Winograd & \(\nmpshape{9}{13}{15}\)  & 1156 & \((5,4)(7,6)(7,8)\) & 1145 \\
Winograd & \(\nmpshape{9}{14}{15}\)  & 1237 & \((4,5)(7,7)(8,7)\) & 1226 \\
\(\nmpshape{2}{2}{3}\) & \(\nmpshape{5}{5}{12}\) & 235 & \((2,3)(2,3)(6,5,\mathbf{1})\) & 232 \\
\bottomrule
\end{tabular}
\end{table}

The last row is the most imbalanced win we found in the sweep
(\(3 \le n \le m \le p \le 16\)): for \(\nmpshape{5}{5}{12}\) under base
\(\nmpshape{2}{2}{3}\) the optimal split of the size-12 axis is
\((6,5,\mathbf{1})\) --- it \emph{peels a width-1 strip}. The products that
touch the unit block compute \(\nmpshape{\cdot}{\cdot}{1}\) sub-shapes at
naive-product rank (here \(R = 4, 6, 6, 9\)), costing almost nothing, while
the remaining products route through the cheap \(\nmpshape{3}{3}{6}{=}40\)
corner --- the unbalanced-extraction mechanism in its purest form.

So the practical rule is not ``prefer balance'' but ``\emph{search} the split,
balance-first'': seed the branch-and-bound with the balanced allocation (an
excellent, usually-optimal incumbent) and let an admissible lower bound prune
the rest. When a known upper bound for the target exists --- typically the
Kronecker rank of a factorisation \(\nmpshape{n}{m}{p} =
\nmpshape{n_1}{m_1}{p_1}\otimes\nmpshape{n_2}{m_2}{p_2}\) --- it is used as the
incumbent: if the base's root lower bound already meets or exceeds it, the base
cannot improve and the allocation sweep is dropped outright.

\subsection{Open questions}
\label{ssec:balance-open}

Two extrapolations are untested here. (i) \emph{Scale}: the catalog thins above
max-dim \(16\), so Table~\ref{tab:omega-imbalance} cannot yet say whether the
average \(\omega\)-vs-imbalance slope persists, flattens, or reverses for large
\(N\); the asymptotic \(\tau\)-theorem constructions suggest the unbalanced tail
should only get richer. (ii) \emph{Rectangular targets}: the metric above
folds rectangular shapes into a single imbalance scalar \(p-n\); a genuinely
rectangular sweep (independent control of \(m-n\) and \(p-m\)) may expose
direction-dependent structure the cubic-centred view misses.
